\section{Учёт токенов и~денежных затрат}\label{sec:analysis}
\subsection{Из чего складывается расход токенов}
Сначала посчитаем расход токенов для одной задачи. Длину её исходного контекста обозначим через $L$. Один цикл планирования, подготовки кандидатов (вариантов программы для решения задачи) и~рецензирования будем называть \emph{раундом}. Зафиксируем число основных раундов $n\ge1$ и~предположим, что оба метода выполняют их полностью. Если возможна остановка до завершения этих раундов, приведённую ниже оценку нужно проверить с~учётом фактического числа раундов.

Если после основных раундов понадобились исправления, все токены дополнительных вызовов, включая повторную передачу контекста, входят в~остаточный расход. Для последовательной процедуры он обозначается $R_{\mathrm{MAS}}$, для интроспективной ветви гибридной процедуры --- $R_{\mathrm H}$.

\Needspace{10\baselineskip}
В~сравниваемом варианте Classical-MAS один раунд содержит $m\ge1$ последовательных отдельных вызовов, каждый из которых выполняет одну роль. Каждый получает весь исходный контекст длины $L$. В~интроспективной ветви Hybrid-INoT на раунд приходится один интроспективный вызов. Длину подготовленного общего контекста в~его $i$-м вызове обозначим через $L_i$. Полный расход токенов равен сумме этих входов и~остальных затрат:
\begin{formulakey}
\begin{align}
T_{\mathrm{MAS}}&=nmL+R_{\mathrm{MAS}},\label{eq:ledger-mas}\\
T_{\mathrm{H}}&=T_{\mathrm{cmp}}+\sum_{i=1}^{n}L_i+R_{\mathrm{H}}.\label{eq:ledger-h}
\end{align}
\begin{symbols}
\sym{T_{\mathrm{MAS}},\,T_{\mathrm H}}{все входные и~выходные токены последовательной процедуры и~интроспективной ветви соответственно}
\sym{n,\ m}{число основных раундов и~число отдельных вызовов в~одном последовательном раунде}
\sym{L,\ nmL}{длина исходного контекста в~токенах и~число токенов всех передач этого контекста в~последовательных раундах}
\sym{i,\ L_i}{номер раунда и~длина подготовленного контекста для его интроспективного вызова}
\sym{T_{\mathrm{cmp}}}{токены подготовки и~сжатия контекста в~интроспективной ветви}
\sym{R_{\mathrm{MAS}},\,R_{\mathrm H}}{остальные токены соответствующего способа. Сюда входят инструкции, ответы, повторно переданные планы, модельные проверки и~все затраты на дополнительные попытки}
\end{symbols}
\end{formulakey}
Каждый токен учитывается один раз. Оба равенства относятся к~фиксированному способу вызова.

Ответы двух процедур и~их длины могут различаться. Эти различия учитываются в~$R_{\mathrm{MAS}}$ и~$R_{\mathrm{H}}$. Если отдельные роли получают только часть контекста, первое слагаемое в~\eqref{eq:ledger-mas} заменяется суммой фактически переданных длин. Таким образом, расчёт зависит от состава каждого вызова.

\subsection{Достаточное условие экономии}
Теперь рассмотрим семейство задач с~разной длиной контекста $L$. Для получения общей оценки зададим ограничения на рост расходов. Применимость этих ограничений к~конкретным моделям и~задачам требует проверки.

\Needspace{14\baselineskip}
\begin{assumption}\label{as:bounds}
Существуют постоянные $\kappa\in(0,1]$, $\gamma,\beta,c_0,c_1,d\ge0$, не зависящие от $L$, такие, что для всех задач рассматриваемого семейства выполнены ограничения:
\begin{formulakey}
\begin{align}
\E[T_{\mathrm{cmp}}]&\le\gamma L+c_0,\\
\E[L_i]&\le\kappa L+d\quad(i=1,\ldots,n),\\
\E[R_{\mathrm{H}}]-\E[R_{\mathrm{MAS}}]&\le\beta L+c_1.
\end{align}
\begin{symbols}
\sym{\E[\cdot]}{среднее по повторным запускам одной процедуры при фиксированной задаче}
\sym{L,\ i,\ n}{длина исходного контекста в~токенах, номер раунда и~число основных раундов}
\sym{T_{\mathrm{cmp}},\,\gamma,\,c_0}{расход подготовки контекста, наклон верхней границы этого расхода и~постоянная добавка в~токенах}
\sym{L_i,\,\kappa,\,d}{длина подготовленного контекста в~раунде~$i$, верхняя доля исходной длины и~добавка в~токенах на раунд, например для диагностики}
\sym{R_{\mathrm H},\,R_{\mathrm{MAS}}}{остаточные токены интроспективной и~последовательной процедур соответственно}
\sym{\beta,\,c_1}{наклон верхней границы разности остаточных расходов и~её постоянная добавка в~токенах}
\end{symbols}
\end{formulakey}
Все средние конечны. Число раундов и~число вызовов в~последовательном раунде фиксированы.
\end{assumption}

Коэффициенты ограничивают рост расходов для выбранного семейства задач. При увеличении длины контекста их значения остаются постоянными.

\Needspace{11\baselineskip}
\begin{proposition}[Условная экономия токенов]\label{prop:savings}
При предположении~\ref{as:bounds} положим
\begin{formulakey}
\begin{equation}
a=n(m-\kappa)-\gamma-\beta,\qquad b=c_0+nd+c_1,
\end{equation}
\begin{symbols}
\sym{a}{наклон нижней границы экономии, в~токенах экономии на токен исходного контекста}
\sym{b}{сумма постоянных накладных расходов в~токенах}
\sym{n,\ m}{число основных раундов и~число последовательных вызовов в~каждом}
\sym{\kappa,\,\gamma,\,\beta}{верхние коэффициенты для длины подготовленного контекста, затрат на его подготовку и~разности остаточных затрат соответственно}
\sym{c_0,\ d,\ c_1}{постоянные добавки на подготовку, на контекст каждого раунда и~на разность остальных затрат соответственно}
\end{symbols}
\end{formulakey}
\begin{formulakey}
Если $a>0$, то для любой рассматриваемой задачи с~$L>\tau_{\mathrm{bound}}=b/a$ выполняется
\nopagebreak
\begin{equation}
\E[T_{\mathrm{MAS}}]-\E[T_{\mathrm{H}}]\ge aL-b>0.
\label{eq:savings}
\end{equation}
\begin{symbols}
\sym{\E[T_{\mathrm{MAS}}]-\E[T_{\mathrm H}]}{средний расход последовательной процедуры за вычетом среднего расхода интроспективной ветви, в~токенах}
\sym{L,\ aL-b}{длина исходного контекста и~нижняя граница экономии в~токенах}
\sym{a,\,b,\,\tau_{\mathrm{bound}}}{прирост нижней границы на токен, постоянные расходы и~их отношение $b/a$, которое должна превысить длина контекста}
\end{symbols}
\end{formulakey}
\end{proposition}
\Needspace{11\baselineskip}
\begin{proof}
Из среднего расхода последовательной процедуры в~\eqref{eq:ledger-mas} вычтем средний расход интроспективной ветви в~\eqref{eq:ledger-h}. У~первой процедуры на исходный контекст затрачено $nmL$ токенов. Чтобы получить нижнюю границу разности, из этой величины вычитаем три верхние оценки из предположения~\ref{as:bounds}: на подготовку контекста, на его передачу в~$n$ раундах и~на разность прочих затрат. После приведения подобных членов коэффициент при длине контекста равен $a$, а~постоянный вычет равен $b$. Отсюда следует~\eqref{eq:savings}. Для сложения средних независимость затрат между раундами не требуется.
\end{proof}

Условие означает, что экономия на повторной передаче контекста должна расти быстрее дополнительных расходов на сжатие, ответы и~повторы. Порог $\tau_{\mathrm{bound}}$ получен из выбранных верхних оценок. Поэтому фактическая точка равенства затрат может находиться в~другом месте. При $a\le0$ знак разности затрат остаётся неопределённым в~рамках этого доказательства.

\noindent\textbf{Числовой пример.} Возьмём один раунд ($n=1$) и~три последовательных вызова ($m=3$). В~интроспективном вызове передаётся весь контекст ($\kappa=1$). Отдельной подготовки и~добавки к~длине контекста нет ($\gamma=c_0=d=0$). Допустим, что разница остальных затрат ограничена половиной исходной длины ($\beta=0{,}5$) и~постоянной добавкой $c_1=1\,000$ токенов. Подстановка даёт
\begin{formulakey}
\[
\begin{aligned}
 a &= 1(3-1)-0-0{,}5 = 1{,}5,\\
 b &= 0+1\cdot 0+1\,000=1\,000,\\
 b/a &=1\,000/1{,}5=666\tfrac{2}{3}.
\end{aligned}
\]
\begin{symbols}
\sym{a}{прирост нижней границы экономии на один токен контекста}
\sym{b}{постоянные накладные расходы в~токенах}
\sym{b/a}{порог длины исходного контекста в~токенах}
\end{symbols}
\end{formulakey}
Длина должна быть строго больше порога. Наименьшая целая длина с~гарантированной экономией в~этом примере составляет 667 токенов. При длине 4\,000 токенов нижняя граница экономии равна $1{,}5\cdot4\,000-1\,000=5\,000$ токенов. Параметры примера иллюстрируют расчёт и~заданы условно.

\subsection{Экономия с~учётом изменения качества}
Зафиксируем задачу $z$ и~настройки обеих процедур. Пусть $Q_{\mathrm{H}}$ и~$Q_{\mathrm{MAS}}$ --- вероятности успешной итоговой проверки $q$ (раздел~\ref{sec:metrics}) при случайном запуске каждой процедуры. Её скрытые тесты могут отличаться от проверки $\mathcal{V}_z$, используемой для исправлений во время работы. Если допускается потеря качества до $\delta$ процентных пунктов относительно Classical-MAS, критерий совместного учёта расхода и~качества имеет вид
\begin{formulakey}
\begin{equation}
\E[T_{\mathrm{H}}]<\E[T_{\mathrm{MAS}}],\qquad 100(Q_{\mathrm{H}}-Q_{\mathrm{MAS}})\ge -\delta.
\label{eq:joint-criterion}
\end{equation}
\begin{symbols}
\sym{\E[T_{\mathrm H}],\,\E[T_{\mathrm{MAS}}]}{средний расход в~токенах при интроспективном и~последовательном способах}
\sym{Q_{\mathrm H},\,Q_{\mathrm{MAS}}}{вероятности успеха тех же способов на независимой проверке}
\sym{100(Q_{\mathrm H}-Q_{\mathrm{MAS}}),\,\delta}{разность долей успеха и~заранее допустимое снижение, обе величины в~процентных пунктах}
\end{symbols}
\end{formulakey}
Первое неравенство следует из утверждения~\ref{prop:savings} при выполнении его предположений. Для второго неравенства нужны данные независимой проверки. При $\delta=0$ критерий требует, чтобы вероятность успеха оставалась как минимум на уровне базового способа.

Для совокупности задач обе части усредняются по одному и~тому же распределению задач. Теоретическая оценка переносится на это среднее, если её предположения и~условие на длину контекста выполняются для всех задач рассматриваемого семейства, а~средние затраты конечны.

Для формальной проверки не меньшей эффективности границу $\delta$ и~уровень доверия задают до анализа данных. Нижняя односторонняя доверительная граница разности должна быть выше $-\delta$. В~серии на 1\,000 задачах рассчитаны разность долей и~её двусторонний интервал (раздел~\ref{subsec:scale}). Протокол этой серии предусматривал сравнение долей с~нулевой разностью и~отдельный анализ расхода. Допустимое снижение качества в~протоколе этой серии не задавалось.

Применение условия экономии также требует проверить, укладываются ли число повторов и~длина сообщений об ошибках в~принятые границы расходов.

\subsection{Верхняя граница полного расхода}
Верхнюю границу полного расхода можно получить сложением лимитов отдельных вызовов. Обозначим через $T_{\mathrm{prep}}^{\max}$ предельный расход подготовки. Для одной итерации зададим пределы $P_{\max}$ входных и~$O_{\max}$ выходных токенов основной модели. Суммарный расход дополнительных языковых моделей в~этой итерации ограничим величиной $S_{\max}$. Тогда
\begin{formulakey}
\begin{equation}
T_{\mathrm{H}}\le T_{\mathrm{prep}}^{\max}+I(P_{\max}+O_{\max}+S_{\max}).
\label{eq:worst}
\end{equation}
\begin{symbols}
\sym{T_{\mathrm H},\,T_{\mathrm{prep}}^{\max}}{полный расход интроспективной ветви и~лимит подготовки в~токенах}
\sym{I}{предельное число итераций}
\sym{P_{\max},\,O_{\max},\,S_{\max}}{лимиты на вход, на ответ основной модели и~на совокупный расход вспомогательных моделей за одну итерацию, в~токенах}
\end{symbols}
\end{formulakey}
Величина $I$ ограничивает число итераций. Контекст, инструкции и~переданная история входят в~$P_{\max}$. Тексты всех $K$ кандидатов и~замечания рецензента уже учтены в~общем пределе ответа $O_{\max}$. Повторное умножение на $K$ завысило бы оценку. Каждая новая попытка занимает одну из $I$ итераций. Поэтому граница сохраняется и~в~случае, когда все попытки неуспешны.

Перед каждым вызовом система резервирует его максимальный расход из оставшегося бюджета. Если такой контроль охватывает и~подготовку, выполняется $T_{\mathrm{H}}\le B$. Для этого сервис должен технически обеспечивать заданные пределы. Скрытые токены рассуждения учитываются по счётчику сервиса. Если он неполон, соответствующий расход отмечается как частично наблюдаемый.

\subsection{Денежная стоимость}
Для расчёта денежной стоимости каждый вид токенов умножается на его цену. Пронумеруем все вызовы моделей за одну задачу индексом $j=1,\ldots,J_{\mathrm{all}}$. Число входных токенов без скидки обозначим через $p_j^{\mathrm{unc}}$, число входных токенов со скидкой за кэширование через $p_j^{\mathrm{cache}}$, а~число оплачиваемых выходных токенов через $r_j$. Кэширование позволяет повторно использовать ранее обработанный вход по правилам сервиса. Соответствующие цены за один токен обозначим через $c_j^{\mathrm{unc}}$, $c_j^{\mathrm{cache}}$, $c_j^{\mathrm{out}}$. Стоимость всех обращений равна
\begin{formulakey}
\begin{equation}
C_{\mathrm{inf}}=\sum_{j=1}^{J_{\mathrm{all}}}
\left(c_j^{\mathrm{unc}}p_j^{\mathrm{unc}}+
c_j^{\mathrm{cache}}p_j^{\mathrm{cache}}+c_j^{\mathrm{out}}r_j\right).
\label{eq:cost}
\end{equation}
\begin{symbols}
\sym{C_{\mathrm{inf}},\,j,\,J_{\mathrm{all}}}{стоимость всех обращений, номер вызова и~число вызовов для одной задачи}
\sym{p_j^{\mathrm{unc}},\,p_j^{\mathrm{cache}},\,r_j}{обычные входные, кэшированные входные и~оплачиваемые выходные токены вызова~$j$}
\sym{c_j^{\mathrm{unc}},\,c_j^{\mathrm{cache}},\,c_j^{\mathrm{out}}}{цены одного токена соответствующих трёх видов в~одной валюте}
\end{symbols}
\end{formulakey}
Величина $C_{\mathrm{inf}}$ выражает стоимость обращений к~моделям. Тариф за миллион токенов переводится в~цену одного токена делением на $10^6$. Для каждого расчёта фиксируются дата и~версия тарифа. Токены рассуждения учитываются один раз в~той категории, по которой они оплачиваются. Стоимость внешних вычислений рассчитывается отдельно.
